% FILE: chapters/ch06_s1_time_series.tex
\chapter{Specialized Areas}\label{ch:specialized_areas}
\section{Time-Series and Sensor Data}\label{sec:time_series}

In conventional machine learning tasks, such as predicting house prices or classifying loan applications, we assume that data samples are \textbf{independent and identically distributed} (often abbreviated as i.i.d.). We can shuffle the rows of the table into any random order without losing any information.

In the Internet of Things (IoT) and edge computing, this assumption is completely false. Sensor readings do not exist in isolation; they arrive as an ordered, continuous chronological stream. The temperature recorded by a sensor at 2:00:01 PM is intimately bound to the temperature recorded at 2:00:00 PM. The sequence of measurements, the speed of change, and the time elapsed between events carry the real physical signal.

This chronological sequence of numerical measurements is called a \textbf{time series}. In this section, we study how to analyze, process, and model temporal data. We examine trend, seasonality, and stationarity, explore classical moving average filters, learn how to slice continuous streams into machine learning windows, compare 1D-CNNs and LSTMs for edge deployment, and evaluate forecasting models without leaking future information.

% ==============================================================================
% 1. TREND, SEASONALITY, AND STATIONARITY
% ==============================================================================
\subsection{Components of Time Series: Trend, Seasonality, and Stationarity}

\begin{keyidea}[Time-Series Decomposition]
Any physical sensor time series $y_t$ can be decomposed into three underlying components:
\begin{enumerate}
    \item \textbf{Trend ($T_t$)}: The long-term gradual increase or decrease over extended time.
    \item \textbf{Seasonality ($S_t$)}: Repeating, cyclical patterns that recur at known, fixed intervals (e.g., daily solar cycles, weekly factory shifts).
    \item \textbf{Residual / Noise ($R_t$)}: Random, irregular fluctuations and electrical measurement noise.
\end{enumerate}
\begin{equation}
\text{Additive Model: } y_t = T_t + S_t + R_t\label{eq:additive_decomp}
\end{equation}
\end{keyidea}

\paragraph{Why It Exists: The Real-World Problem}
Imagine monitoring an agricultural soil moisture sensor over six months. If soil moisture drops from $45\%$ to $20\%$, is the irrigation system failing, or is it simply mid-afternoon during a hot summer week? By decomposing the signal, we separate the daily diurnal evaporation cycle (seasonality) from a true mechanical pipe blockage (trend).

\begin{analogy}[Ocean Tides and Wind Waves]
Stand on a beach and observe the ocean water level:
\begin{itemize}
    \item \textbf{Seasonality}: The predictable, rhythmic rise and fall of the tide every 12 hours caused by the gravitational pull of the moon.
    \item \textbf{Trend}: The slow, multi-year rise in global sea levels.
    \item \textbf{Noise}: The chaotic, choppy surface ripples generated by random gusts of wind.
\end{itemize}
\end{analogy}

\paragraph{The Concept of Stationarity}
Most statistical and machine learning algorithms assume that the underlying data distribution is \textbf{stationary}.
\begin{keyidea}[Stationarity]
A time series is \textbf{strictly stationary} if its statistical properties—specifically its \textbf{mean}, \textbf{variance}, and \textbf{autocorrelation}—do not change over time.
\begin{itemize}
    \item \textbf{Constant Mean}: The series does not drift upward or downward over time.
    \item \textbf{Constant Variance}: The amplitude of fluctuations remains stable.
    \item \textbf{Constant Autocorrelation}: The relationship between $y_t$ and $y_{t-k}$ depends only on the lag $k$, not on the calendar date $t$.
\end{itemize}
\end{keyidea}

If a sensor series has a strong rising trend, its mean is continuously shifting. A model trained on data with a mean of $20^\circ\text{C}$ will fail completely when the physical system drifts into an operating mean of $60^\circ\text{C}$.

\paragraph{Transforming Non-Stationary Data: Differencing}
To transform non-stationary data into a stationary series, we compute the difference between consecutive time steps. The \textbf{first-order difference} ($\Delta y_t$) is:
\begin{equation}
\Delta y_t = y_t - y_{t-1}\label{eq:first_diff}
\end{equation}
Instead of modeling the absolute level, we model the \textit{change per unit time}.

\paragraph{Hand Calculation Example}
Consider four consecutive hourly temperature readings from a greenhouse experiencing morning heating:
\begin{equation*}
y = [18.0, 20.0, 23.0, 27.0]^\circ\text{C}
\end{equation*}
This raw series is clearly non-stationary because the mean is steadily rising. Let us compute the first-order differenced series $\Delta y$:
\begin{align*}
\Delta y_1 &= y_1 - y_0 = 20.0 - 18.0 = +2.0^\circ\text{C} \\
\Delta y_2 &= y_2 - y_1 = 23.0 - 20.0 = +3.0^\circ\text{C} \\
\Delta y_3 &= y_3 - y_2 = 27.0 - 23.0 = +4.0^\circ\text{C}
\end{align*}
Differencing removed the baseline offset of $18.0^\circ\text{C}$, converting the raw sensor values into the rate of temperature rise ($+2, +3, +4$).

\paragraph{Python Code}
\begin{lstlisting}[language=Python]
import pandas as pd
import numpy as np

# Simulated sensor series: baseline trend + daily cycle + noise
np.random.seed(42)
time_index = pd.date_range("2026-01-01", periods=100, freq="h")
trend = np.linspace(10, 30, 100)
seasonality = 5 * np.sin(2 * np.pi * np.arange(100) / 24)
noise = np.random.normal(0, 0.5, 100)

raw_series = pd.Series(trend + seasonality + noise, index=time_index, name="temp_c")

# First-order differencing
diff_series = raw_series.diff().dropna()

print("--- Raw Non-Stationary Series ---")
print(f"Mean:     {raw_series.mean():.2f} C")
print(f"Variance: {raw_series.var():.2f}")

print("\n--- First-Order Differenced Series ---")
print(f"Mean:     {diff_series.mean():.2f} C/hour")
print(f"Variance: {diff_series.var():.2f}")
\end{lstlisting}

\begin{lstlisting}[style=outputstyle]
--- Raw Non-Stationary Series ---
Mean:     20.03 C
Variance: 46.54

--- First-Order Differenced Series ---
Mean:     0.20 C/hour
Variance: 1.05
\end{lstlisting}

\paragraph{Line-by-Line Explanation}
\begin{itemize}
    \item Lines 6--9: We construct an additive synthetic time series containing a steady rising slope from $10$ to $30^\circ\text{C}$, a 24-hour sinusoidal wave, and Gaussian noise.
    \item Line 14: \code{raw\_series.diff()} computes $y_t - y_{t-1}$. The first entry becomes \code{NaN} because there is no prior value, so we remove it with \code{.dropna()}.
    \item Lines 17--23: Notice how the variance drops from $46.54$ down to $1.05$. Differencing stabilized the variance and reduced the drifting mean to a steady $+0.20^\circ\text{C/hour}$.
\end{itemize}

\begin{takeaways}
\begin{itemize}
    \item Time series data violates the i.i.d. assumption; sequential temporal order is mandatory.
    \item Decomposition separates long-term trends and repeating seasonal cycles from random noise.
    \item Differencing stabilizes drifting means, transforming non-stationary data into stationary inputs.
\end{itemize}
\end{takeaways}

% ==============================================================================
% 2. MOVING AVERAGES AND ARIMA
% ==============================================================================
\subsection{Classical Forecasting: Moving Averages and ARIMA}

\begin{keyidea}[Classical Forecasting Models]
Classical statistical models predict future values based entirely on mathematical patterns observed in past values:
\begin{itemize}
    \item \textbf{Moving Average (MA)}: A smoothing filter that removes short-term noise by averaging recent observations.
    \item \textbf{ARIMA (AutoRegressive Integrated Moving Average)}: A statistical model combining lagged observations (AR), differencing (I), and lagged forecast errors (MA).
\end{itemize}
\end{keyidea}

\paragraph{Why It Exists: The Real-World Problem}
Physical sensors connected to microcontrollers are noisy. An analog temperature sensor on an ESP8266 fluctuates by $\pm 0.8^\circ\text{C}$ every second due to electrical switching noise on the power supply rail. If an actuator triggers a cooling fan the instant temperature crosses $30.0^\circ\text{C}$, the fan will chatter on and off several times per second, rapidly burning out the relay coil. Moving average filters smooth out this noise.

\paragraph{Exponential Moving Average (EMA)}
A Simple Moving Average (SMA) requires storing the last $K$ readings in a memory buffer. On a microcontroller with tiny RAM, an \textbf{Exponential Moving Average (EMA)} is far superior: it smooths data recursively using \textbf{only a single scalar of memory}!
\begin{equation}
S_t = \alpha y_t + (1 - \alpha) S_{t-1}\label{eq:ema_formula}
\end{equation}
\begin{itemize}
    \item $y_t$ is the new raw incoming measurement at time $t$.
    \item $S_t$ is the updated smooth estimate.
    \item $S_{t-1}$ is the previous smoothed value stored in memory.
    \item $\alpha \in (0, 1]$ is the \textbf{smoothing factor}. A higher $\alpha$ reacts faster to new readings; a lower $\alpha$ filters out more noise.
\end{itemize}

\paragraph{Hand Calculation Example: Exponential Moving Average}
Let an ESP32 initialize its smoothed temperature estimate to $S_0 = 20.0^\circ\text{C}$. We configure a smoothing factor of $\alpha = 0.2$. Two consecutive noisy readings arrive: $y_1 = 30.0^\circ\text{C}$ (a brief electrical noise spike) followed by $y_2 = 21.0^\circ\text{C}$:
\begin{enumerate}
    \item \textbf{Update on First Reading ($y_1 = 30.0$)}:
    \begin{equation*}
    S_1 = \alpha y_1 + (1 - \alpha) S_0 = (0.2 \times 30.0) + (0.8 \times 20.0) = 6.0 + 16.0 = \mathbf{22.0^\circ\text{C}}
    \end{equation*}
    Notice that the sudden $+10^\circ\text{C}$ spike was damped down to only $+2^\circ\text{C}$!
    \item \textbf{Update on Second Reading ($y_2 = 21.0$)}:
    \begin{equation*}
    S_2 = \alpha y_2 + (1 - \alpha) S_1 = (0.2 \times 21.0) + (0.8 \times 22.0) = 4.2 + 17.6 = \mathbf{21.8^\circ\text{C}}
    \end{equation*}
    The filter smoothly adapted to the true physical state without allocating any array buffers in RAM.
\end{enumerate}

\paragraph{The ARIMA$(p, d, q)$ Model}
For formal forecasting, the ARIMA model combines three parameters:
\begin{itemize}
    \item $\mathbf{p}$ (\textbf{AutoRegressive order}): Number of past lagged observations used to predict the current value:
    \begin{equation*}
    y_t = c + \phi_1 y_{t-1} + \phi_2 y_{t-2} + \dots + \phi_p y_{t-p} + \epsilon_t
    \end{equation*}
    \item $\mathbf{d}$ (\textbf{Degree of differencing}): Number of times the raw series must be differenced to achieve stationarity.
    \item $\mathbf{q}$ (\textbf{Moving Average order}): Number of past forecast error residuals used in the prediction:
    \begin{equation*}
    y_t = \mu + \epsilon_t + \theta_1 \epsilon_{t-1} + \dots + \theta_q \epsilon_{t-q}
    \end{equation*}
\end{itemize}

\paragraph{Python Code}
\begin{lstlisting}[language=Python]
import pandas as pd
import numpy as np
from statsmodels.tsa.arima.model import ARIMA

# Simulated battery discharge time series (hourly % capacity)
np.random.seed(42)
steps = 40
true_capacity = 100.0 - 1.5 * np.arange(steps) + np.random.normal(0, 0.4, steps)
series = pd.Series(true_capacity)

# Split into train (35 hours) and test (5 hours)
train, test = series[:35], series[35:]

# Fit an ARIMA(1, 1, 0) model: 1 lag, 1 difference, 0 error terms
model = ARIMA(train, order=(1, 1, 0))
fitted_model = model.fit()

# Forecast 5 steps into the future
forecast = fitted_model.forecast(steps=5)

print("--- ARIMA(1, 1, 0) Forecast vs Actual ---")
comparison = pd.DataFrame({"Actual": test.values, "Forecast": forecast.values.round(2)})
print(comparison)
\end{lstlisting}

\begin{lstlisting}[style=outputstyle]
--- ARIMA(1, 1, 0) Forecast vs Actual ---
   Actual  Forecast
0   47.78     47.66
1   46.61     46.17
2   44.33     44.68
3   43.25     43.19
4   41.97     41.70
\end{lstlisting}

\begin{iotbox}[Two-Line Real-Time EMA Filter on ESP32 in C++]
On an ESP32 or Arduino microcontroller reading an ADC pin, you can implement an EMA low-pass filter in two lines of C++ code:
\begin{lstlisting}[language=C++]
float filtered_temp = 20.0;  // Preserved in global memory
float alpha = 0.15;

void loop() {
  float raw_temp = analogRead(34) * (3.3 / 4095.0) * 100.0;
  filtered_temp = (alpha * raw_temp) + ((1.0 - alpha) * filtered_temp);
  delay(100);
}
\end{lstlisting}
This eliminates high-frequency noise spikes completely while consuming only 4 bytes of RAM.
\end{iotbox}

\begin{takeaways}
\begin{itemize}
    \item Exponential Moving Average (EMA) filters noise recursively with $O(1)$ memory, perfect for microcontrollers.
    \item The smoothing factor $\alpha$ balances noise suppression against sensor responsiveness.
    \item ARIMA$(p, d, q)$ captures autoregressive lags and forecast errors to project future trends.
\end{itemize}
\end{takeaways}

% ==============================================================================
% 3. WINDOWING SENSOR STREAMS
% ==============================================================================
\subsection{Windowing Sensor Streams for Machine Learning}

\begin{keyidea}[The Sliding Window Technique]
Standard supervised machine learning models require a 2D matrix of input features $\mathbf{X}$ and target labels $\mathbf{y}$. 
The \textbf{sliding window technique} transforms a continuous, one-dimensional time stream into structured feature vectors by sweeping a fixed-length window across the time axis.
\end{keyidea}

\paragraph{Why It Exists: The Real-World Problem}
An accelerometer on a factory motor streams 100 vibrations per second. A neural network cannot accept an endless, indefinite stream of numbers. To classify whether the motor is healthy, imbalanced, or experiencing bearing failure, we slice the stream into discrete windows (e.g., 50 consecutive samples representing 0.5 seconds of motion) and feed each window to the classifier.

\paragraph{The Parameters of a Sliding Window}
\begin{itemize}
    \item \textbf{Window Size ($W$)}: The number of historical time steps contained in each input sample.
    \item \textbf{Stride / Step Size ($S$)}: The number of time steps the window advances forward for the next sample. If $S < W$, windows \textbf{overlap}.
    \item \textbf{Prediction Horizon ($H$)}: How many steps into the future the model must predict.
\end{itemize}

\begin{figure}[htbp]
\centering
\begin{tikzpicture}[scale=0.9, >=Stealth]
    % Time axis
    \draw[->, thick] (0, 0) -- (11, 0) node[right, font=\small\sffamily] {Time $t$};
    \foreach \t in {0,1,2,3,4,5,6,7,8,9,10} {
        \draw (\t, 0.1) -- (\t, -0.1) node[below, font=\tiny\sffamily] {$t_{\t}$};
    }

    % Window 1 (W=3, S=1, H=1)
    \draw[thick, keyideablue, fill=keyideablue!10, rounded corners=1mm] (0, 0.6) rectangle (3, 1.3);
    \node[keyideablue, font=\footnotesize\bfseries\sffamily] at (1.5, 0.95) {Window 1: $[t_0, t_1, t_2]$};
    \draw[->, thick, warningred] (3, 0.95) -- (4, 0.95) node[right, font=\footnotesize\bfseries\sffamily] {Target $y = t_3$};

    % Window 2 (Stride = 1)
    \draw[thick, analogygreen, fill=analogygreen!10, rounded corners=1mm] (1, 1.7) rectangle (4, 2.4);
    \node[analogygreen, font=\footnotesize\bfseries\sffamily] at (2.5, 2.05) {Window 2: $[t_1, t_2, t_3]$};
    \draw[->, thick, warningred] (4, 2.05) -- (5, 2.05) node[right, font=\footnotesize\bfseries\sffamily] {Target $y = t_4$};

    % Window 3 (Stride = 1)
    \draw[thick, iotpurple, fill=iotpurple!10, rounded corners=1mm] (2, 2.8) rectangle (5, 3.5);
    \node[iotpurple, font=\footnotesize\bfseries\sffamily] at (3.5, 3.15) {Window 3: $[t_2, t_3, t_4]$};
    \draw[->, thick, warningred] (5, 3.15) -- (6, 3.15) node[right, font=\footnotesize\bfseries\sffamily] {Target $y = t_5$};
\end{tikzpicture}
\caption{The sliding window mechanism: slicing a 1D stream into supervised input-output pairs $(W=3, S=1, H=1)$.}\label{fig:sliding_window}
\end{figure}

\paragraph{Hand Calculation Example}
Given a raw stream of six temperature measurements:
\begin{equation*}
\text{Stream} = [15.0, 16.0, 18.0, 21.0, 25.0, 30.0]
\end{equation*}
Let us extract all supervised samples using window size $W = 3$, stride $S = 1$, and horizon $H = 1$:
\begin{align*}
\text{Sample 1: } \mathbf{x}_1 &= [15.0, 16.0, 18.0], \quad y_1 = 21.0 \\
\text{Sample 2: } \mathbf{x}_2 &= [16.0, 18.0, 21.0], \quad y_2 = 25.0 \\
\text{Sample 3: } \mathbf{x}_3 &= [18.0, 21.0, 25.0], \quad y_3 = 30.0
\end{align*}
The 1D stream of 6 numbers was transformed into a tabular dataset of 3 samples, each having 3 features and 1 target label.

\paragraph{Python Code}
\begin{lstlisting}[language=Python]
import numpy as np

def create_sliding_windows(series: np.ndarray, window_size: int, horizon: int = 1, stride: int = 1):
    """Transforms a 1D time series array into 2D supervised feature pairs (X, y)."""
    X, y = [], []
    total_steps = len(series)
    
    for i in range(0, total_steps - window_size - horizon + 1, stride):
        window = series[i : i + window_size]
        target = series[i + window_size + horizon - 1]
        X.append(window)
        y.append(target)
        
    return np.array(X), np.array(y)

# Example sensor stream
raw_sensor_stream = np.array([15.0, 16.0, 18.0, 21.0, 25.0, 30.0])

X_windows, y_targets = create_sliding_windows(
    series=raw_sensor_stream, 
    window_size=3, 
    horizon=1, 
    stride=1
)

print("Generated Feature Matrix X (Shape: ", X_windows.shape, "):\n", X_windows)
print("\nGenerated Target Vector y  (Shape: ", y_targets.shape, "):\n", y_targets)
\end{lstlisting}

\begin{lstlisting}[style=outputstyle]
Generated Feature Matrix X (Shape:  (3, 3) ):
 [[15. 16. 18.]
 [16. 18. 21.]
 [18. 21. 25.]]

Generated Target Vector y  (Shape:  (3,) ):
 [21. 25. 30.]
\end{lstlisting}

\begin{takeaways}
\begin{itemize}
    \item Sliding windows convert continuous time streams into supervised training matrices $\mathbf{X}$ and $\mathbf{y}$.
    \item Window size $W$ dictates how much past context the model considers.
    \item Overlapping windows ($S < W$) augment training sample count while preserving local patterns.
\end{itemize}
\end{takeaways}

% ==============================================================================
% 4. DEEP LEARNING FOR TIME SERIES: LSTM AND 1D-CNN
% ==============================================================================
\subsection{Deep Learning for Time Series: LSTM and 1D-CNN}

\begin{keyidea}[Temporal Deep Learning Architectures]
When temporal patterns are too complex for linear regression, we deploy two deep learning architectures:
\begin{itemize}
    \item \textbf{1D Convolutional Neural Network (1D-CNN)}: Slides 1D filter kernels across time to detect local features (such as rapid frequency spikes or sudden impacts).
    \item \textbf{Long Short-Term Memory (LSTM)}: A recurrent architecture that maintains a continuous internal memory cell across time, regulating information flow through input, forget, and output gates.
\end{itemize}
\end{keyidea}

\paragraph{Why It Exists: The Real-World Problem}
In Human Activity Recognition (HAR) using smartwatch accelerometers, motion patterns (walking, running, falling) involve rhythmic cycles that evolve over several seconds. A standard dense MLP treats every time index as an independent feature and cannot easily recognize that a vibration spike occurring at sample 10 is identical to a spike occurring at sample 35. 1D-CNNs provide \textbf{translation invariance across time}.

\paragraph{1D Convolution: Slicing with Filter Kernels}
In a 1D-CNN, a filter kernel $\mathbf{k}$ of size $K$ slides along the sequence $\mathbf{x}$, computing dot products at each step:
\begin{equation}
(x * k)[t] = \sum_{m=0}^{K-1} x_{t+m} k_m\label{eq:conv1d}
\end{equation}

\paragraph{Hand Calculation Example: 1D Edge/Difference Kernel}
Let us slide a difference-detection filter kernel $\mathbf{k} = [-1.0, 1.0]$ across a sensor signal $\mathbf{x} = [2.0, 5.0, 1.0, 4.0]$:
\begin{enumerate}
    \item \textbf{Step $t=0$}:
    \begin{equation*}
    (x * k)[0] = (-1.0 \times 2.0) + (1.0 \times 5.0) = -2.0 + 5.0 = \mathbf{+3.0}
    \end{equation*}
    \item \textbf{Step $t=1$}:
    \begin{equation*}
    (x * k)[1] = (-1.0 \times 5.0) + (1.0 \times 1.0) = -5.0 + 1.0 = \mathbf{-4.0}
    \end{equation*}
    \item \textbf{Step $t=2$}:
    \begin{equation*}
    (x * k)[2] = (-1.0 \times 1.0) + (1.0 \times 4.0) = -1.0 + 4.0 = \mathbf{+3.0}
    \end{equation*}
    Output Feature Map: $[+3.0, -4.0, +3.0]$.
\end{enumerate}
The kernel detected the sharp rise ($+3$), steep fall ($-4$), and subsequent rebound ($+3$).

\paragraph{The LSTM Cell: Gated Memory}
Standard Recurrent Neural Networks (RNNs) suffer from vanishing gradients when sequences exceed 20 or 30 time steps. LSTMs solve this using three gates:
\begin{enumerate}
    \item \textbf{Forget Gate ($f_t$)}: $\sigma(\mathbf{W}_f [\mathbf{h}_{t-1}, \mathbf{x}_t] + b_f)$ decides what percentage of old memory to wipe.
    \item \textbf{Input Gate ($i_t$)}: Decides which new information to write into the cell state.
    \item \textbf{Output Gate ($o_t$)}: Decides what portion of internal memory to output as hidden state $\mathbf{h}_t$.
\end{enumerate}

\paragraph{Python Code with PyTorch}
\begin{lstlisting}[language=Python]
import torch
import torch.nn as nn

# Input batch: 4 samples, 10 time steps, 3 sensor channels (e.g., ax, ay, az)
# Shape required by PyTorch: (batch_size, channels, sequence_length) for Conv1d
sample_batch_cnn = torch.randn(4, 3, 10)

# 1. 1D Convolutional Layer for temporal pattern extraction
conv1d = nn.Conv1d(in_channels=3, out_channels=8, kernel_size=3)
out_cnn = conv1d(sample_batch_cnn)

print(f"1D-CNN Input Shape:  {sample_batch_cnn.shape} (batch, channels, time)")
print(f"1D-CNN Output Shape: {out_cnn.shape} (batch, 8 filters, time - 2)")

# 2. LSTM Layer for sequence modeling
# Shape required by PyTorch LSTM: (batch_size, sequence_length, features)
sample_batch_lstm = torch.randn(4, 10, 3)
lstm = nn.LSTM(input_size=3, hidden_size=16, batch_first=True)
out_lstm, (hn, cn) = lstm(sample_batch_lstm)

print(f"\nLSTM Input Shape:   {sample_batch_lstm.shape} (batch, time, features)")
print(f"LSTM Output Shape:  {out_lstm.shape} (batch, time, hidden_units)")
\end{lstlisting}

\begin{lstlisting}[style=outputstyle]
1D-CNN Input Shape:  torch.Size([4, 3, 10]) (batch, channels, time)
1D-CNN Output Shape: torch.Size([4, 8, 8]) (batch, 8 filters, time - 2)

LSTM Input Shape:   torch.Size([4, 10, 3]) (batch, time, features)
LSTM Output Shape:  torch.Size([4, 10, 16]) (batch, time, hidden_units)
\end{lstlisting}

\begin{iotbox}[Why 1D-CNNs Dominate TinyML Over LSTMs]
On microcontrollers like the ESP32 or ARM Cortex-M4, \textbf{1D-CNNs are vastly preferred over LSTMs}:
\begin{itemize}
    \item \textbf{Compute Efficiency}: 1D convolutions can be computed in parallel and compile directly to ARM CMSIS-NN SIMD (Single Instruction Multiple Data) assembly instructions.
    \item \textbf{Memory Footprint}: LSTMs must maintain internal state buffers for every step in the recurrence chain, consuming precious SRAM. 1D-CNNs run with lower memory overhead and execute up to $5\times$ faster on edge silicon.
\end{itemize}
\end{iotbox}

\begin{takeaways}
\begin{itemize}
    \item 1D-CNNs slide kernels across time to extract translation-invariant local vibration patterns.
    \item LSTMs maintain internal memory cells through gated circuits, preventing vanishing gradients over long horizons.
    \item 1D-CNNs are significantly faster and more memory-efficient on edge microcontrollers than recurrent models.
\end{itemize}
\end{takeaways}

% ==============================================================================
% 5. FORECASTING EVALUATION
% ==============================================================================
\subsection{Forecasting Evaluation: Metrics and Temporal Validation}

\begin{keyidea}[Temporal Validation and Forecast Metrics]
In time-series modeling, \textbf{you must never evaluate models using random train/test splits}. Random shuffling allows the model to train on future data to predict the past, causing catastrophic \textbf{temporal data leakage}.
\end{keyidea}

\paragraph{Why It Exists: The Real-World Problem}
Suppose you train a stock or weather predictor by randomly shuffling days into training and test sets. The model uses Wednesday's temperature to predict Tuesday's temperature. When evaluated in cross-validation, it reports an impossible $99.9\%$ accuracy. But when deployed in the real world on Thursday morning, it fails completely because it does not have access to Friday's data!

\paragraph{Temporal Cross-Validation (Walk-Forward Validation)}
In time series, we use \textbf{Walk-Forward Cross-Validation} (also called rolling-origin validation):
\begin{itemize}
    \item Fold 1: Train on Days 1--10 $\implies$ Test on Day 11.
    \item Fold 2: Train on Days 1--11 $\implies$ Test on Day 12.
    \item Fold 3: Train on Days 1--12 $\implies$ Test on Day 13.
\end{itemize}
The training set only ever contains historical data that occurred strictly \textbf{before} the test period.

\paragraph{Standard Forecast Evaluation Metrics}
For $N$ time steps with true values $y_t$ and predictions $\hat{y}_t$:
\begin{enumerate}
    \item \textbf{Mean Absolute Error (MAE)}:
    \begin{equation}
    \text{MAE} = \frac{1}{N} \sum_{t=1}^N |y_t - \hat{y}_t|\label{eq:mae_ts}
    \end{equation}
    Easy to interpret because it is expressed in the exact physical units of the sensor.

    \item \textbf{Root Mean Squared Error (RMSE)}:
    \begin{equation}
    \text{RMSE} = \sqrt{\frac{1}{N} \sum_{t=1}^N (y_t - \hat{y}_t)^2}\label{eq:rmse_ts}
    \end{equation}
    Heavily penalizes occasional large spikes or catastrophic failures.

    \item \textbf{Mean Absolute Percentage Error (MAPE)}:
    \begin{equation}
    \text{MAPE} = \frac{100\%}{N} \sum_{t=1}^N \left| \frac{y_t - \hat{y}_t}{y_t} \right|\label{eq:mape_ts}
    \end{equation}
    Expresses error as a relative percentage of actual value.
\end{enumerate}

\paragraph{Hand Calculation Example}
A solar forecast model makes two predictions:
\begin{itemize}
    \item Day 1: Actual $y_1 = 20.0\text{ kW}$, Predicted $\hat{y}_1 = 22.0\text{ kW}$
    \item Day 2: Actual $y_2 = 25.0\text{ kW}$, Predicted $\hat{y}_2 = 21.0\text{ kW}$
\end{itemize}
Let us calculate all three metrics by hand:
\begin{enumerate}
    \item \textbf{Errors}:
    $e_1 = |20.0 - 22.0| = 2.0\text{ kW}$, $e_2 = |25.0 - 21.0| = 4.0\text{ kW}$.
    \item \textbf{MAE}:
    \begin{equation*}
    \text{MAE} = \frac{2.0 + 4.0}{2} = \frac{6.0}{2} = \mathbf{3.0\text{ kW}}
    \end{equation*}
    \item \textbf{RMSE}:
    \begin{equation*}
    \text{RMSE} = \sqrt{\frac{2.0^2 + 4.0^2}{2}} = \sqrt{\frac{4.0 + 16.0}{2}} = \sqrt{\frac{20.0}{2}} = \sqrt{10.0} \approx \mathbf{3.1623\text{ kW}}
    \end{equation*}
    \item \textbf{MAPE}:
    \begin{align*}
    \text{PE}_1 &= \frac{2.0}{20.0} = 0.10 \quad (10.0\%) \\
    \text{PE}_2 &= \frac{4.0}{25.0} = 0.16 \quad (16.0\%) \\
    \text{MAPE} &= \frac{10.0\% + 16.0\%}{2} = \frac{26.0\%}{2} = \mathbf{13.0\%}
    \end{align*}
\end{enumerate}

\paragraph{Python Code}
\begin{lstlisting}[language=Python]
import numpy as np
from sklearn.model_selection import TimeSeriesSplit

# Time-series cross-validation demonstration
data = np.array([10, 20, 30, 40, 50, 60, 70, 80])

tscv = TimeSeriesSplit(n_splits=3)
print("--- TimeSeriesSplit Indices (Train vs Test) ---")
for fold, (train_idx, test_idx) in enumerate(tscv.split(data)):
    print(f"Fold {fold+1}: Train Indices {train_idx} -> Test Indices {test_idx}")
\end{lstlisting}

\begin{lstlisting}[style=outputstyle]
--- TimeSeriesSplit Indices (Train vs Test) ---
Fold 1: Train Indices [0 1 2 3 4] -> Test Indices [5]
Fold 2: Train Indices [0 1 2 3 4 5] -> Test Indices [6]
Fold 3: Train Indices [0 1 2 3 4 5 6] -> Test Indices [7]
\end{lstlisting}

\begin{takeaways}
\begin{itemize}
    \item Always use forward-chaining temporal validation to avoid leaking future data into training.
    \item MAE reflects average error magnitude in physical units; RMSE penalizes severe outlier spikes.
    \item MAPE normalizes error as a percentage, but divides by zero if actual values equal $0$.
\end{itemize}
\end{takeaways}

% ==============================================================================
% 6. CHECK YOURSELF
% ==============================================================================
\subsection{Check Yourself}

Test your understanding of time-series analysis and sensor stream processing:

\begin{enumerate}
    \item \textbf{Stationarity \& Differencing}: Why is a raw temperature sensor stream exhibiting a steady $+0.5^\circ\text{C}$ rise per hour considered non-stationary? What does first-order differencing do to this stream?

    \item \textbf{Exponential Moving Average Calculation}: An ESP32 environmental monitor has a current filtered humidity estimate $S_0 = 50.0\%$. Two new readings arrive from a DHT22 sensor: $y_1 = 60.0\%$ followed by $y_2 = 40.0\%$. Given smoothing factor $\alpha = 0.25$, compute $S_1$ and $S_2$ by hand.

    \item \textbf{Sliding Window Sizing}: An industrial vibration accelerometer streams data at $100\text{ Hz}$ for 10 seconds ($1{,}000$ readings total). You configure a sliding window with window size $W = 100$ samples, stride $S = 50$ samples, and horizon $H = 1$. Calculate the total number of window samples $N$ produced.

    \item \textbf{Architecture Comparison for TinyML}: Why do embedded machine learning engineers favor 1D-CNNs over LSTMs when deploying vibration anomaly detection models to microcontrollers like the ESP32?

    \item \textbf{Evaluation Pitfalls}: Why does evaluating a time-series model with Scikit-Learn's standard \code{train_test_split(shuffle=True)} produce unrealistically optimistic test scores?
\end{enumerate}

% ==============================================================================
% ANSWERS
% ==============================================================================
\subsection*{Answers}

\begin{enumerate}
    \item \textbf{Stationarity \& Differencing}:
    \begin{itemize}
        \item A stream with a steady rise of $+0.5^\circ\text{C/hour}$ is non-stationary because its **mean is continuously drifting over time** (violating the constant mean condition).
        \item First-order differencing ($\Delta y_t = y_t - y_{t-1}$) subtracts consecutive readings, transforming the drifting raw temperatures into a constant series: $[0.5, 0.5, 0.5, \dots]$. The differenced series has a constant mean of $0.5$ and zero trend, making it stationary and suitable for modeling.
    \end{itemize}

    \item \textbf{Exponential Moving Average Calculation}:
    Using $S_t = \alpha y_t + (1 - \alpha) S_{t-1}$ with $\alpha = 0.25$:
    \begin{align*}
    S_1 &= (0.25 \times 60.0) + (0.75 \times 50.0) = 15.0 + 37.5 = \mathbf{52.5\%} \\
    S_2 &= (0.25 \times 40.0) + (0.75 \times 52.5) = 10.0 + 39.375 = \mathbf{49.375\%}
    \end{align*}

    \item \textbf{Sliding Window Sizing}:
    Total readings $T = 1{,}000$. Window length $W = 100$, stride $S = 50$, horizon $H = 1$.
    Using the window count formula:
    \begin{equation*}
    N = \left\lfloor \frac{T - W - H}{S} \right\rfloor + 1 = \left\lfloor \frac{1000 - 100 - 1}{50} \right\rfloor + 1 = \left\lfloor \frac{899}{50} \right\rfloor + 1 = 17 + 1 = \mathbf{18 \text{ windows}}
    \end{equation*}

    \item \textbf{Architecture Comparison for TinyML}:
    \begin{itemize}
        \item **Memory Footprint**: LSTMs require retaining cell state memory buffers across sequential recurrence steps, which exhausts limited microcontroller SRAM. 1D-CNNs require no recurrence buffers.
        \item **Compute Speed**: 1D-CNNs perform feedforward convolutions that can be parallelized and mapped directly to hardware SIMD assembly instructions (e.g., ARM CMSIS-NN), executing up to $5\times$ faster with lower energy consumption.
    \end{itemize}

    \item \textbf{Evaluation Pitfalls}:
    Shuffling data randomly breaks the temporal timeline, scattering future observations into the training split and past observations into the test split. The model effectively learns to ``predict yesterday by peeking at tomorrow'' (temporal data leakage). When deployed in production, future data is unavailable, causing the model to fail despite high offline test metrics.
\end{enumerate}
